\documentclass[11pt,a4paper]{article}

\usepackage[utf8]{inputenc}
\usepackage[T1]{fontenc}
\usepackage[spanish,es-noshorthands]{babel}
\usepackage{lmodern}
\usepackage{microtype}
\usepackage[a4paper,margin=2.5cm]{geometry}
\usepackage{amsmath}
\usepackage{booktabs}
\usepackage{array}
\usepackage{enumitem}
\usepackage{xcolor}
\usepackage{tikz}
\usetikzlibrary{positioning,arrows.meta,shapes.geometric}
\usepackage{listings}
\usepackage[hidelinks]{hyperref}

\setlength{\emergencystretch}{2.5em}

\definecolor{codebg}{RGB}{246,246,246}
\lstset{
  basicstyle=\ttfamily\footnotesize,
  backgroundcolor=\color{codebg},
  frame=single,
  rulecolor=\color{gray!40},
  breaklines=true,
  columns=fullflexible,
  keepspaces=true,
  showstringspaces=false,
  commentstyle=\color{gray!70!black},
  keywordstyle=\color{blue!60!black}\bfseries
}

\title{\textbf{Seminario de Arquitecturas de Software}\\[2pt]
\large Respuesta a la Pregunta 3:\\
Aplicación de DDD y \emph{Bounded Context} en un Dominio Complejo}
\author{}
\date{}

\begin{document}
\maketitle

\section*{3. Aplicación de DDD y \emph{Bounded Context} en un dominio complejo}

\subsection*{Planteamiento}

Una aseguradora desea construir un sistema para gestionar pólizas de automóviles que
interactúe con tres subsistemas con lenguajes y preocupaciones propias: el departamento
de siniestros, el sistema de facturación de primas y la aplicación móvil de los ajustadores
que evalúan daños en el lugar del accidente. El término ``póliza'' no significa lo mismo en
cada uno de ellos: para un actuario es un contrato con cláusulas y coberturas, para
facturación es una fuente de ingresos y de cuotas, y para un ajustador es apenas un número
de referencia que identifica el siniestro que debe inspeccionar.

Intentar un único modelo de \texttt{Poliza} que satisfaga a todos esos actores conduce al
problema clásico que Evans (2003) denomina \emph{modelo único para todo el negocio}: un
modelo sobrecargado, con campos y comportamientos que solo interesan a un área, acoplado a
bases de datos compartidas y cada vez más difícil de modificar. La respuesta de
\emph{Domain-Driven Design} (DDD) no es construir un modelo más grande, sino dividir el
dominio en varios \emph{contextos delimitados} (\emph{Bounded Contexts}), cada uno con su
propio \emph{lenguaje ubicuo} y su propio modelo, y definir entre ellos relaciones
explícitas y contratos de integración (Evans, 2003; Richards y Ford, 2020).

\subsection*{1. División del sistema en contextos delimitados}

La primera decisión de diseño es delimitar las fronteras del modelo, no las de la
organización ni las de la interfaz de usuario. Cada contexto debe agrupar un conjunto
coherente de conceptos y reglas que un mismo equipo puede mantener y que usan un mismo
lenguaje. Para este sistema se proponen los siguientes contextos:

\begin{table}[h]
\centering
\footnotesize
\setlength{\tabcolsep}{4pt}
\renewcommand{\arraystretch}{1.25}
\begin{tabular}{@{}p{2.7cm}p{3.9cm}p{5.1cm}p{2.9cm}@{}}
\toprule
\textbf{Contexto} & \textbf{Lenguaje ubicuo} & \textbf{Responsabilidad / modelo} & \textbf{Tipo (DDD)} \\
\midrule
\textbf{Suscripción y Pólizas} &
Póliza, cláusula, cobertura, vigencia, endoso, asegurado, vehículo &
Contrato de seguro: emisión, modificaciones, cancelación, reglas de cobertura y versión del contrato. Es la fuente de verdad del contrato. &
Núcleo (\emph{core domain}) \\
\addlinespace
\textbf{Siniestros} &
Siniestro, reclamación, ocurrencia, beneficiario, dictamen, cobertura aplicable &
Registro y tramitación del siniestro; determina qué coberturas aplican según la póliza vigente al momento del hecho. &
Núcleo \\
\addlinespace
\textbf{Facturación de Primas} &
Prima, cuota, recibo, mora, nota de crédito, conciliación &
Generación y cobro de las obligaciones económicas derivadas de la póliza; no interpreta cláusulas, solo montos y vencimientos. &
Soporte \\
\addlinespace
\textbf{Ajustes en Campo} &
Ajuste, inspección, daño, evidencia, peritaje, número de referencia &
Asignación, ejecución y cierre de inspecciones de daños; trabaja con datos mínimos del siniestro y del vehículo. &
Soporte \\
\addlinespace
\textbf{Catálogo y Tarifas} &
Producto, plan, tarifa, factor de riesgo, cobertura comercial &
Definición de productos y precios que consumen Suscripción y Facturación. &
Genérico / soporte \\
\addlinespace
\textbf{Identidad y Notificaciones} &
Usuario, ajustador, asegurado, canal, plantilla &
Autenticación, autorización y envío de comunicaciones. &
Genérico \\
\bottomrule
\end{tabular}
\caption{Contextos delimitados propuestos y su lenguaje ubicuo.}
\end{table}

La delimitación sigue el criterio de Evans (2003): una frontera de contexto es una frontera
de modelo y de lenguaje. Nótese que ``Póliza'' aparece en varios contextos, pero con
significados y ciclos de vida diferentes; por eso existen varias clases \texttt{Poliza},
una por contexto, y no una clase compartida. Esto es intencional: compartir el modelo
sería precisamente la fuente de contaminación que se quiere evitar.

\subsection*{2. Cómo evitar que el modelo de ``Póliza'' se contamine}

La contaminación se evita con cuatro mecanismos complementarios.

\begin{enumerate}[leftmargin=1.4em]
\item \textbf{Un modelo por contexto y propiedad exclusiva de los datos.} Cada contexto
posee su esquema de base de datos y solo ese contexto escribe en él. Siniestros y
Facturación no acceden a las tablas de Suscripción; obtienen la información por un
contrato de integración. No existen \emph{joins} ni claves foráneas entre contextos, ni
entidades JPA/EF compartidas entre módulos.

\item \textbf{Referencias por identidad, no por objeto.} En Siniestros, la póliza se
representa mediante un objeto de valor \texttt{PolicyId} y, cuando se requiere el detalle,
mediante una \emph{fotografía} (\emph{snapshot}) inmutable de las cláusulas y coberturas
vigentes en la fecha del siniestro. Así, una modificación posterior de la póliza (un
endoso) no altera retroactivamente un siniestro ya dictaminado, y el modelo de Siniestros
no necesita conocer la estructura interna del contrato.

\item \textbf{Capa anticorrupción (\emph{Anti-Corruption Layer}, ACL).} Todo dato que
provenga de otro contexto pasa por un traductor que convierte el modelo externo en el
modelo local. La ACL protege el lenguaje ubicuo: si el sistema de facturación cambia sus
campos o su semántica, solo se modifica el traductor, no el dominio de Siniestros
(Evans, 2003; Richards y Ford, 2020).

\item \textbf{Fronteras verificadas automáticamente.} Las dependencias entre contextos se
restringen por paquetes/namespaces y se validan en integración continua con pruebas de
arquitectura (por ejemplo, ArchUnit en la JVM o reglas equivalentes en .NET). Un contexto
no puede importar clases de dominio de otro; solo puede usar sus contratos de integración.
\end{enumerate}

En consecuencia, el modelo de Póliza del contexto de Suscripción conserva únicamente
conceptos contractuales (cláusulas, coberturas, vigencias, endosos). Facturación modela
una \texttt{CuentaDeCobro} con cuotas y recibos derivados de eventos, y Ajustes modela una
\texttt{AsignacionDeAjuste} con un número de referencia, sin heredar ni replicar el
contrato completo.

\subsection*{3. Comunicación desacoplada entre contextos}

Los contextos se integran mediante \textbf{eventos de dominio publicados como eventos de
integración} y mediante \textbf{contratos explícitos} (interfaces de consulta o APIs),
nunca mediante una base de datos compartida. Los patrones de \emph{context mapping} de
Evans (2003) permiten nombrar cada relación:

\begin{itemize}[leftmargin=1.4em]
\item \textbf{Suscripción como \emph{Open Host Service} con \emph{Published Language}}: publica
un contrato estable (por ejemplo, un esquema de eventos versionado) que otros contextos
consumen sin conocer su modelo interno.
\item \textbf{Siniestros y Facturación como consumidores con ACL}: cada uno traduce esos
eventos a su propio modelo. La relación es de tipo \emph{customer--supplier}: el
consumidor define sus necesidades y el proveedor se compromete a mantener el contrato.
\item \textbf{Catálogo y Tarifas como proveedor conforme} (\emph{conformist}) o con ACL,
según la rigidez del contrato: si el catálogo es corporativo e inmutable, los demás se
adaptan a él; si no, se interpone una ACL.
\item \textbf{Identidad y Notificaciones como servicio genérico} consumido por todos,
con relaciones de tipo \emph{separate ways} respecto del núcleo.
\end{itemize}

\begin{sloppypar}
Los eventos típicos son \texttt{PolicyIssued}, \texttt{PolicyEndorsed},
\texttt{PolicyCancelled}, \texttt{ClaimRegistered}, \texttt{ClaimApproved} y
\texttt{PaymentReceived}. La comunicación es \textbf{asíncrona} a través de un bus de
mensajes, lo que desacopla en el tiempo y en el espacio a emisores y consumidores y evita
que una caída de Facturación detenga la atención de un siniestro. Para no perder eventos
se aplica el patrón \emph{outbox} transaccional, y para no duplicar efectos, consumidores
idempotentes. Los procesos que abarcan varios contextos (por ejemplo, aprobar un siniestro
y generar la indemnización en Facturación) se coordinan con una \emph{saga} o
\emph{process manager}, aceptando consistencia eventual (Richards y Ford, 2020).
\end{sloppypar}

\begin{figure}[h]
\centering
\begin{tikzpicture}[
  font=\small, scale=0.92, transform shape,
  ctx/.style={draw, rounded corners, align=center, minimum width=3.1cm,
              minimum height=1.15cm, fill=blue!4},
  core/.style={ctx, fill=orange!12, thick},
  rel/.style={-{Latex[length=2mm]}, thick, gray!70!black},
  lbl/.style={font=\scriptsize, fill=white, inner sep=1pt}
]
\node[core] (pol) {\textbf{Suscripción y Pólizas}\\{\scriptsize Núcleo · OHS/PL}};
\node[ctx, right=2.6cm of pol] (sin) {\textbf{Siniestros}\\{\scriptsize Núcleo · ACL}};
\node[ctx, below=1.7cm of sin] (fac) {\textbf{Facturación de Primas}\\{\scriptsize Soporte · ACL}};
\node[ctx, below=1.7cm of pol] (aju) {\textbf{Ajustes en Campo}\\{\scriptsize Soporte · ACL}};
\node[ctx, left=1.6cm of pol, yshift=1.0cm] (cat) {\textbf{Catálogo y Tarifas}\\{\scriptsize Genérico}};
\node[ctx, left=1.6cm of pol, yshift=-1.0cm] (idn) {\textbf{Identidad y Notif.}\\{\scriptsize Genérico}};

\draw[rel] (pol) -- node[lbl, above]{\texttt{PolicyIssued}} (sin);
\draw[rel] (pol.south) -- node[lbl, left]{\texttt{PolicyEndorsed}} (aju.north);
\draw[rel] (sin.south) -- node[lbl, right]{\texttt{ClaimRegistered}} (fac.north);
\draw[rel] (sin) -- node[lbl, above right]{\texttt{ClaimApproved}} (fac);
\draw[rel, dashed] (cat) -- (pol);
\draw[rel, dashed] (idn) -- (pol);
\end{tikzpicture}
\caption{Mapa de contextos: contratos de integración y eventos principales.}
\end{figure}

\subsection*{4. Relación con una arquitectura orientada al dominio (Onion/Clean)}

El diseño anterior es independiente del estilo de despliegue, pero encaja naturalmente con
la Arquitectura Limpia y la Arquitectura Cebolla: \textbf{cada contexto delimitado se
implementa como un núcleo de dominio aislado}, y las integraciones son detalles
reemplazables (Martin, 2017; Palermo, 2008).

\begin{itemize}[leftmargin=1.4em]
\item \textbf{Centro (dominio):} entidades, objetos de valor, agregados, eventos de dominio
y reglas invariantes del contexto. No conoce frameworks, bases de datos ni HTTP.
\item \textbf{Aplicación:} casos de uso y \emph{puertos} de entrada y de salida. La Regla de
Dependencia garantiza que el código fuente solo apunte hacia el centro.
\item \textbf{Infraestructura:} adaptadores que implementan los puertos (repositorios,
mensajería, clientes HTTP). Aquí viven las ACL y los traductores entre modelos.
\item \textbf{Interfaces:} controladores REST, consumidores de eventos y la aplicación móvil
del ajustador, como adaptadores de entrada.
\end{itemize}

La ``Guía de Arquitectura N-Capas orientada al Dominio'' de De la Torre Llorente et al.
(2010) describe este mismo mecanismo con el nombre de \textbf{Adaptador de Servicio}: un
componente que traduce el contrato de un servicio externo al contrato que espera la capa
de aplicación, aislando al dominio de la tecnología del proveedor. En términos de
Arquitectura Limpia, el Adaptador de Servicio es un \emph{adaptador de salida} que
implementa un \emph{puerto} definido en el núcleo. Por ejemplo, cuando Siniestros necesita
conocer la póliza, el dominio define el puerto y la infraestructura aporta el adaptador:

\begin{lstlisting}[language=Caml, caption={Puerto en el núcleo y Adaptador de Servicio en infraestructura.}, label={lst:puerto}]
// Nucleo (capa de dominio/aplicacion): el puerto pertenece al contexto
public interface PolicyQueryPort {
    PolicySnapshot findSnapshotBy(PolicyId id, LocalDate occurrenceDate);
}

// Infraestructura: Adaptador de Servicio + Anti-Corruption Layer
public sealed class PolicyServiceAdapter implements PolicyQueryPort {
    private final PolicyApiClient client;      // detalle tecnologico
    private final PolicyTranslator translator; // ACL

    @Override
    public PolicySnapshot findSnapshotBy(
            PolicyId id, LocalDate date) {
        PolicyDto dto = client.getPolicy(id.value(), date);
        return translator.toSnapshot(dto); // sin PolicyDto en el dominio
    }
}
\end{lstlisting}

El dominio de Siniestros depende de \texttt{PolicyQueryPort}, no de
\texttt{PolicyApiClient}; la inyección de dependencias se resuelve en el
\emph{composition root} de infraestructura. El mismo esquema aplica a la publicación y
consumo de eventos: el dominio declara \texttt{DomainEventPublisher} y la infraestructura
lo implementa sobre el bus de mensajes. Así, la Arquitectura Orientada al Dominio no solo
organiza las capas de cada contexto, sino que hace explícitos los límites entre contextos:
cada frontera coincide con un puerto y con un contrato versionado.

\subsection*{5. Consecuencias y conclusión}

La división propuesta produce un sistema en el que (i) cada área mantiene su propio
lenguaje y su propio ritmo de cambio; (ii) el modelo de Póliza del núcleo permanece
protegido por identidades, fotografías y ACL; (iii) los contextos se comunican por eventos
y contratos versionados, sin bases de datos compartidas; y (iv) cada contexto puede
implementarse como un núcleo Onion/Clean con adaptadores de servicio intercambiables. El
costo es asumir consistencia eventual y una mayor complejidad operativa, que se justifica
en un dominio complejo donde la autonomía de los equipos y la evolucionabilidad del
negocio son requisitos de primer orden (Evans, 2003; Richards y Ford, 2020).

\section*{Referencias}

\begin{thebibliography}{9}
\setlength{\itemsep}{2pt}

\bibitem[Evans, 2003]{evans2003}
Evans, E. (2003). \emph{Domain-Driven Design: Tackling Complexity in the Heart of
Software}. Addison-Wesley.

\bibitem[Cockburn, 2005]{cockburn2005}
Cockburn, A. (2005). \emph{Hexagonal Architecture (Ports and Adapters)}.
\url{https://alistair.cockburn.us/hexagonal-architecture/}

\bibitem[De la Torre Llorente et al., 2010]{guia2010}
De la Torre Llorente, C., Zorrilla Castro, U., Ramos Barroso, M. A., y Calvarro Nelson, J.
(2010). \emph{Guía de Arquitectura N-Capas orientada al Dominio con .NET 4.0}
(Borrador). Microsoft Ibérica.

\bibitem[Martin, 2017]{martin2017}
Martin, R. C. (2017). \emph{Clean Architecture: A Craftsman's Guide to Software Structure
and Design}. Prentice Hall.

\bibitem[Palermo, 2008]{palermo2008}
Palermo, J. (2008). \emph{The Onion Architecture}.
\url{https://jeffreypalermo.com/2008/07/the-onion-architecture-part-1/}

\bibitem[Richards y Ford, 2020]{richards2020}
Richards, M., y Ford, N. (2020). \emph{Fundamentals of Software Architecture: An
Engineering Approach}. O'Reilly Media.

\end{thebibliography}

\end{document}
